\chapter{Introduction}

Quel est le problème à résoudre ? Pourquoi y a-t-il un problème ? Quelle est la solution envisagée ? Pourquoi est-ce que la solution envisagée est meilleure que l'état de l'art ?

En bref, on vous demande d'introduire votre travail, l'idée de départ et les objectifs attendus. Le lecteur qui découvrirait votre projet au travers de cette introduction devrait ainsi être capable d'en comprendre le cadre, l'idée générale et les aboutissants du projet.

\section{Contexte}

Enjeux, objetifs généraux et but

\section{Problématique}

La problématique constitue le cœur de l'introduction et représente le point de départ de votre réflexion. Elle doit être clairement définie afin de présenter de manière explicite le problème que vous souhaitez résoudre. Une problématique bien formulée permet au lecteur de comprendre les enjeux du projet, le contexte dans lequel il s'inscrit et les raisons pour lesquelles cette question mérite d'être traitée.

Il ne s'agit pas seulement d'exposer un sujet, mais d'identifier une difficulté spécifique, un manque ou un besoin auquel votre projet apporte une réponse. Vous devez également souligner la pertinence de votre démarche : pourquoi ce problème est-il important ? Quelle valeur ajoutée votre travail est-il susceptible d'apporter ? En clarifiant ces points, vous démontrerez la légitimité et l'intérêt de votre projet.

\section{Objectifs}

Objectifs précis SMART (Spécifique, Mesurable, Atteignable, Réaliste, Temporellement défini)

\section{Méthodologie}

Vous devez justifier chacune de vos décisions : pourquoi avez-vous opté pour une telle approche plutôt qu'une autre ? Quels sont les avantages de cette méthode par rapport aux alternatives disponibles ? Cette section permet également de présenter les critères de sélection des outils, les méthodes d'analyse utilisées, ainsi que le cadre expérimental mis en place.
